\section{The product of independent degree scales}\label{sec:product}

Continue with the cone and slice constructed in Section~\ref{sec:cones}.
For $1\leq j\leq n$, let $d_j$ be the first positive normalized degree
at which the family of \emph{all} homogeneous functions up to that
degree has transcendence rank at least $j$.  These ranks can also be
attained using finite-group eigenfunctions, by choosing character bases
in each homogeneous piece.

\begin{proposition}\label{prop:product}
After passage to a subsequence there is a uniform constant $C$ such that
\begin{equation}\label{prod:product-bound}
 h\,d_1\cdots d_n\leq C,
 \qquad h=\vol(v).
\end{equation}
\end{proposition}

The proof cuts the projective slice by homogeneous pencils and compares
the volumes of successive geometric generic fibres.  The following
estimate supplies the discrepancy bound used in those comparisons.
Its proof in Section~\ref{sec:integral-jump} is independent of
Proposition~\ref{prop:product}, and its index-one hypothesis concerns
the volume form itself.

\begin{proposition}[Integral jump]\label{prop:integral-jump}
Fix a positive dimension.  Consider a sequence of smooth projective
varieties $F$ over algebraically closed fields isomorphic to $\C$, with
projective birational morphisms $F\to F_0$ to normal varieties.  Let
$\eta_F$ be an lc rational volume form, and let $L$ be the pullback of
a big semiample effective $\Q$-Cartier divisor on $F_0$.  Suppose the
supports of $\operatorname{div}(\eta_F)$ and $L$ on $F$ have simple
normal crossings.  Write $l_w=w(L)$ and assume the following conditions.
\begin{enumerate}[label=\textup{(\roman*)}]
\item For every divisorial valuation $w$, the equality
      $A_{\eta_F}(w)=0$ implies $l_w>0$.
\item For every prime divisor $P$ on $F_0$,
      \[
      A_{\eta_F}(P)=1\quad\text{if }l_P=0,
      \qquad
      A_{\eta_F}(P)\leq\epsilon l_P\quad\text{if }l_P>0,
      \]
      where $\epsilon\to0$ along the sequence.
\item For every rational $0<b\leq1/10$,
      $h^0(F,K_F+\lceil bL\rceil)=1$.
\item There is a uniform constant $C$ such that every effective
      $\Q$-divisor $J\sim_{\Q}L$ satisfies
      $w(J)\leq C l_w$ at every divisorial valuation with
      $A_{\eta_F}(w)=0$.
\end{enumerate}
Then, after passage to a subsequence, there is a uniform constant $C'$
such that every such $J$ and every divisorial valuation $w$ satisfy
\begin{equation}\label{prod:integral-jump-bound}
 w(J)\leq C'\bigl(A_{\eta_F}(w)+l_w\bigr).
\end{equation}
It suffices that the hypotheses hold sufficiently far along the sequence.
\end{proposition}

\subsection{Degree tiers and geometric generic components}

By Proposition~\ref{prop:cone-estimates} and
\eqref{cone:complement-section},
\[
 \frac1{2n}\leq d_1\leq N.
\]
Passing to a subsequence, split the ordered degrees into successive
tiers: ratios within a tier are bounded above and below, whereas the
ratio from the end of one tier to the beginning of the next tends to
infinity.  This follows by successively passing to subsequences for the
finitely many adjacent ratios.  Every degree in the first tier is
bounded above and below.

Choose algebraically independent homogeneous functions
$f_1=f_*,f_2,\ldots,f_n$ of normalized degrees $e_1=N,e_2,\ldots,e_n$,
where $e_j$ is comparable with $d_j$.  Choose them so that the functions
in each complete prefix of tiers form a transcendence basis for the
homogeneous functions before the next tier.  To do this, extend the
single bounded-degree element $f_*$ already in the first tier; its
degree may exceed that tier's initial cutoff by a bounded amount, but
it lies below the next tier sufficiently far along the sequence.
Thereafter extend the existing independent family at each tier cutoff.
No equivariance of these functions is required.  If $1,\ldots,m$ is a
complete prefix and $T=d_{m+1}$ starts the next tier, every homogeneous
element of normalized degree strictly less than $T$ is algebraic over
$k(f_1,\ldots,f_m)$, where $k$ is the cone's ground field.

For any prefix $1,\ldots,m$, in particular for a complete prefix of
tiers, define rational functions on $\widehat E$ by
\begin{equation}\label{prod:pencils}
 \phi_j=\frac{f_j}{u^{ae_j}},\qquad 2\leq j\leq m.
\end{equation}
They are independent, since $f_1=u^{aN}$.  Resolve the graph of their
map to $(\mathbb P^1)^{m-1}$, obtaining a smooth projective $W$, and
let $F$ be a connected component of its geometric generic fibre.
Set $q=n-m$.  Let $F_0$ be the corresponding component on the
normalized graph; it is normal by localization and geometric base
change in characteristic zero.  The divisor $D$ pulls back to divisors
$D_W,D_F$, and to an effective $\Q$-Cartier divisor on $F_0$.
The divisor $D_F$ is nef and semiample.  It is also big: the geometric
generic component meets densely the graph over the original domain of
definition, so its map to $\widehat E$ is birational onto its image,
and $D$ is ample there.

Divide $\eta$ by
$d\phi_2\wedge\cdots\wedge d\phi_m$, with a fixed choice of sign, to
obtain the relative volume form $\eta_F$.  It is lc, and
\begin{equation}\label{prod:lc-poles}
 A_{\eta_F}(w)=0\quad\Longrightarrow\quad w(D_F)>0.
\end{equation}
Indeed, on an SNC resolution generic smoothness restricts the divisor
of the form and its discrepancy formula to the geometric generic
fibre.  Off $D$ the form is klt, which gives the implication.  The
algebraically closed fields of these fibres are isomorphic to $\C$:
adjoining finitely many transcendental elements and then taking an
algebraic closure leaves the transcendence degree over
$\overline{\Q}$ equal to that of $\C$.

We may choose one $W$ dominating all prefix graphs, with every pencil
regular and the supports of $\eta,D_W$ SNC.  Apply generic smoothness
also to every stratum.  On every intermediate generic fibre the
restricted support then has normal crossings and its cuts are reduced.
Consequently taking ceilings of rational multiples of $D_W$ commutes
with these restrictions.  Components missing the generic fibre
contribute nothing.  The absolute-to-relative division defining
$\eta_F$ likewise restricts the form divisor over the generic
parameters.

\subsection{Discrepancies on the normal graph}

\begin{lemma}\label{prod:graph-bound}
At every prime divisor $P$ of $F_0$, with its primitive valuation, put
$c_P=P(D_F)$.  Then
\begin{equation}\label{prod:graph-estimate}
 \begin{aligned}
 c_P>0&\quad\Longrightarrow\quad
 A_{\eta_F}(P)\leq
          \left(\sum_{j\leq m}e_j-1\right)c_P,\\
 c_P=0&\quad\Longrightarrow\quad A_{\eta_F}(P)=1.
 \end{aligned}
\end{equation}
This includes the case $m=1$.
\end{lemma}

\begin{proof}
Spread $P$ to a horizontal divisor valuation $P_0$ on the normalized
graph over the rational base.  A finite algebraic extension of the
generic base used to label its component is unramified at this
horizontal divisor.  By separability on the divisor, the order of the
relative differential form is its absolute order before division by
the base differentials.  Thus we can compute using $P_0$.

When $c_P=0$, its generic centre lies off $D$.  The rational map is
regular there, so its normalized graph is isomorphic to that open of
$\widehat E$, where $\eta$ has zero Weil divisor.  This gives discrepancy
one.  Suppose $c_P>0$ and let $S$ be
the centre of $P_0$ on $\widehat E$, of codimension $r\geq1$.
Among the residues of the $\phi_j$ there are $r-1$ algebraically
independent elements over $k(S)$.  To see this, the residue field of
$P_0$ has transcendence degree $n-2$, and is algebraic over the field
generated by $k(S)$ and the base coordinates: normalization of the
graph is finite over the graph.  All these base coordinates are units
at the horizontal divisor, so a transcendence basis can be chosen from
their residues.  The difference of transcendence degrees is $r-1$.

Extend $P_0$ to $k(\widehat E)(u)$ by the Gauss valuation $\nu$ with
$\nu(u)=c_P/a$, and also restrict $\nu$ to the original cone field.
It is nonnegative on homogeneous regular elements: the divisor of a
degree-$I$ section divided by $u^I$ is bounded below by $-IH_u$.
Thus $\nu$ has a centre on the affine cone cover.

The field of this affine centre has transcendence degree $n-r$ and is
algebraically related to $k(S)(\tau)$ over a common subfield of that
transcendence degree.  Here $\tau$ is the residue of a positive-degree
section that is a unit for $\nu$.  More explicitly, choose a very ample
Cartier multiple and a section nonvanishing at $S$; its residue is
transcendental over the residue field of $P_0$ by the Gauss construction.
Ratios of other sections in that multiple recover the dimensions of
$S$.  For every other homogeneous unit, taking powers and ratios with
the chosen section makes its residue algebraic over $k(S)(\tau)$.
Homogeneous decomposition proves the same assertion for residues of
units in the entire affine ring.  The finite map to $\Spec R$ preserves
the dimension of this centre.

Choose $n-r$ regular functions on $\Spec R$ with independent residues
on that centre.  The selected $r-1$ residues of the $\phi_j$, or their
positive powers, remain independent over its centre field, because
$\tau$ is Gauss-transcendental over the residue field of $P_0$.
Together with $f_1$, use the corresponding $r-1$ functions among the
$f_j$.  After taking powers, their degree-zero ratios with $f_1$ give
the chosen powers of the $\phi_j$.

These $n$ cone functions have a logarithmic differential wedge of
discrepancy zero at $\nu$.  In fact monomial combinations with nonzero
determinant give one function of positive order and $n-1$ units with
independent residues.  In characteristic zero their logarithmic wedge
has a simple pole along the divisorial valuation.  Hence their ordinary
differential wedge has discrepancy equal to the sum of their values,
which is at most $\sum_{j\leq m}e_jc_P$.  Dividing that wedge by
$\theta$ gives a regular coefficient on the cone: this holds at every
codimension-one point, and normality extends it.  On the other hand,
\eqref{cone:slice-form} and the logarithmic Gauss formula give
\[
 A_\theta(\nu)=A_{\eta_F}(P)+c_P.
\]
Nonnegativity of the coefficient's value proves
\eqref{prod:graph-estimate}.  If $m=1$, necessarily $r=1$ in this
calculation; the list of $\phi_j$ is empty and the same proof applies.
\end{proof}

\subsection{A single adjoint section before a degree jump}

\begin{lemma}\label{prod:one-section}
Suppose $q>0$ and $T=d_{m+1}$ begins the tier after the chosen complete
prefix.  Uniformly sufficiently far along that jump,
\begin{equation}\label{prod:adjoint-unique}
 h^0\bigl(F,K_F+\lceil bTD_F\rceil\bigr)=1
 \qquad(0<b\leq1/10,\ b\in\Q).
\end{equation}
\end{lemma}

\begin{proof}
The form $\eta_F$ provides a section.  Its poles are simple and lie
over $D_F$, so the positive ceiling supplies at least one copy of each
pole.  Suppose a second section has nonconstant ratio with $\eta_F$.
After extending scalars to the algebraic closure of the rational base
field, regard $F$ as a component of a general smooth complete
intersection of the free pencils on $W$.  Extend the section by zero
on the other components.  Let $M_j$ be the pole divisor on $W$ of
$\phi_j$; it is a member of the corresponding free pencil.

Successive adjunction lifts this section to
\[
 K_W+\lceil bTD_W\rceil+\sum_{j=2}^m M_j.
\]
To give the restriction step explicitly, on an intermediate smooth
component $Y$ write $A=bTD_Y$, let $C$ be the next general member of
its free pencil, and let $R_Y$ be the sum of the remaining pencil
divisors.  Adjunction gives the exact sequence
\[
\begin{split}
0\longrightarrow&\ \cO_Y(K_Y+\lceil A\rceil+R_Y)\\
\longrightarrow&\ \cO_Y(K_Y+\lceil A\rceil+R_Y+C)\\
\longrightarrow&\ \cO_C(K_C+\lceil A|_C\rceil+R_Y|_C)
\longrightarrow0.
\end{split}
\]
The first term has vanishing $H^1$ by Kawamata--Viehweg vanishing
\cite[Section~3.1]{FujinoFundamental}.  Indeed,
$B_Y=\lceil A\rceil-A$ has SNC support and coefficients in $[0,1)$,
and the difference of the displayed integral divisor and $K_Y+B_Y$
is $A+R_Y$, which is big and nef.  Bigness on each intermediate
component follows, as for $F$, from its image on the original domain
of the graph.  All supports were resolved before taking the cuts, and
ceilings commute with these cuts as explained above.  In the finite
base chart, adjunction by the
equations $\phi_j-\lambda_j$ identifies the lifted ratio with the
required ratio of relative forms.  Denote the lift, written as a
rational multiple of $\eta$, by $H'\eta$.

On the base-changed $\widehat E$, the poles of $H'$ are bounded by
\begin{equation}\label{prod:lift-poles}
 \left(bT+N+\sum_{j=2}^m e_j\right)D.
\end{equation}
Indeed, at a prime with coefficient $c>0$ in $D$, the case $m=1$ of
Lemma~\ref{prod:graph-bound} gives
$\ord\eta\leq(N-1)c-1$.  The ceiling excess is at most one, and the
pole divisor of each $\phi_j$ is at most $e_jD$.  Adding these
inequalities proves \eqref{prod:lift-poles}.  At a prime with $c=0$ the
form has order zero and these divisors have no pole.

The selected root field is cyclic Galois over the original cone field,
of degree dividing $aN$.  Split $H'$ into its deck-character parts,
with the enlarged scalar field fixed.  Each part still obeys
\eqref{prod:lift-poles}, because $D$ is invariant; the deck group is
not required to preserve the selected geometric fibre.  For a part of
given character choose an integer $I$ in the cancelling congruence
class with
\[
 bT+N+\sum_{j=2}^m e_j
 \leq \frac Ia
 <bT+N+\sum_{j=2}^m e_j+N+\frac1a.
\]
Multiplication by $u^I$ makes it an invariant homogeneous regular
function on the cone cover.  Regularity follows in codimension one
from the pole inequality and the divisor $H_u$; the vertex has
codimension at least two, so normality completes the check.  This
invariant is therefore in the scalar extension of $R_I$.

The preceding degrees are negligible compared with $T$, and $a\geq1$.
Since $bT\leq T/10$, one index in the sequence sufficiently far along
the jump makes $I/a<T$ for every rational $b$ in the asserted range.
Every ground-field basis element $h\in R_I$ is then algebraic over
$k(f_1,\ldots,f_m)$.  Its quotient $h/u^I$ is algebraic over
$k(\phi_2,\ldots,\phi_m)$.  To check the latter assertion, after
adjoining $u$ the original algebraic relation is a relation over
$k(\phi_2,\ldots,\phi_m)(u)$; but $h/u^I$ belongs to
$k(\widehat E)$ and $u$ is transcendental over that field.  Expanding
a relation in powers of $u$ gives a nonzero relation over the base
field alone.  Such a rational function is constant on each connected
geometric generic component.  This remains true for scalar combinations
of the basis elements.  Thus every deck-character part of $H'$ is
constant on $F$, contradicting the chosen nonconstant ratio.
\end{proof}

\subsection{Jets, image sheaves and clipping on the first tier}

Take the complete bounded first tier, of length $m$; allow $m=n$ if
there is only one tier.  We prove
\begin{equation}\label{prod:initial-volume}
 h\leq C D_F^q,\qquad D_F^q\leq Ch.
\end{equation}
For $q=0$, the degree on the geometric generic component is one.
If $q>0$, we also prove that, for every divisorial valuation $P$ over
$F$ with $A_{\eta_F}(P)=0$, every effective $J\sim_{\Q}D_F$ satisfies
\begin{equation}\label{prod:clipping}
 P(J)\leq Cc_P,\qquad c_P=P(D_F)>0.
\end{equation}

The upper volume bound follows from nef intersections on $W$.
Both $D_W$ and the pencil divisors $M_j$ are nef, and
$e_jD_W-M_j$ is effective.  Therefore
\[
 D_F^q\leq D_W^q\prod_{j=2}^m M_j
 \leq\left(\prod_{j=2}^m e_j\right)D_W^{n-1}
 \leq Ch
\]
by \eqref{cone:volume} and boundedness of the first tier.  If a fibre
has several components, the first intersection is their sum with
positive multiplicities, so it still bounds the chosen component.

For the other direction test the full vector space
$R_{\leq ap}=\bigoplus_{i\leq ap}R_i$, with $p$ large and divisible.
Write one of its elements as $H=\sum_{i\leq ap}h_i^\circ$, where
$h_i^\circ\in R_i$.  On $W\times\mathbb G_m$, with torus coordinate
$s$, it becomes
\begin{equation}\label{prod:test-functions}
 H_s=\sum_{i\leq ap}s^i\frac{h_i^\circ}{u^i}.
\end{equation}
It is a rational section with pole bound $pD_W$.  For every rational
divisorial valuation $w$ centred on this product and every nonzero $H$,
\begin{equation}\label{prod:simultaneous-order}
 w(H_s)\leq Cp\bigl(A_{\eta\wedge d\log s}(w)+w(D_W)\bigr).
\end{equation}
To prove this estimate, add a Gauss parameter $z$ with positive rational
weight $\gamma>w(D_W)/a$, and use the generically finite map to the
cone times the $z$-line given by dilation $s\mapsto sz$.  Every
positive-degree homogeneous cone function now has positive value, so
the valuation is centred at the vertex times zero.  Let $I$ be the
largest occurring degree.  If $I=0$ the assertion is immediate.
Otherwise test the single function
$\sum_i z^{I-i}h_i^\circ$ on the target.  At $z=0$ it is $h_I^\circ$,
whose threshold is at least $1/(2np)$ by \eqref{cone:alpha}.
Inversion of adjunction, with the central divisor added, gives
\[
 w(H_s)+I\gamma
 \leq 2np\bigl(A_{\eta\wedge d\log s}(w)+a\gamma\bigr).
\]
The discrepancy on the right follows by pulling back
$\theta\wedge dz/z$ and using \eqref{cone:slice-form}; the torus
coordinate $s$ is a unit.  Let $\gamma$ decrease to $w(D_W)/a$ and
discard the nonnegative term on the left to obtain
\eqref{prod:simultaneous-order}.  The calculation also works after a
generically finite base change, by restricting valuations and pulling
back the same forms and divisors.

We now give the jet count, retaining the exact section space at a
thickened divisor.  For the divisor test, first make $P$ primitive.
Spread its extraction after a finite generically \'etale base change
of an open of the pencil base, labelling the selected component $F$.
All geometric generic data descend over a finite extension.  Resolve
and shrink so that the family is projective, the required relative
SNC data are smooth, and the chosen components and divisor descend
geometrically integrally.  Include $s$ as an unchanged extra base
parameter.  At the generic point of the specialized divisor on a
sufficiently general closed fibre, take monomial weights one on the
$m$ base parameters, including $s-s_0$, and weight $1/c_P$ on the
equation of the spread divisor.  The relative form has a simple pole
there and no other support component passes through this generic point.
The resulting rational divisorial valuation satisfies
\begin{equation}\label{prod:jet-test-values}
 w(D_W)=1,\qquad A_{\eta\wedge d\log s}(w)=m.
\end{equation}
For the whole-component test use instead the generic point of a
component of a general closed fibre and only the $m$ base parameters,
all of weight one.  Then the divisor value is zero and discrepancy is
$m$.  Commensurability makes both valuations divisorial up to scaling.
The indicated parameters form a regular system of parameters at the
respective generic points.  The nonzero functions being tested remain
nonzero by dominance.

By \eqref{prod:simultaneous-order}, the order of each nonzero test
section as a local section of $\cO(pD_W)$ is at most $C_1p$ for either
valuation.  The constant is uniform in $p$, $P$ and the sufficiently
general closed parameters, although those parameters may be chosen
separately for each $p,P$.

Write $\pi:\mathcal X\to B_0$ for the spreading component family and
$\mathcal P$ for the spread divisor.  In the whole-component test use
$\pi_*\cO_{\mathcal X}(pD_W)$.  In the divisor test use the
\emph{image sheaf}
\[
 \mathcal E_{p,j}=\operatorname{im}\left(
 \pi_*\cO_{\mathcal X}(pD_W)
 \longrightarrow\pi_*\cO_{j\mathcal P}(pD_W)\right),
 \qquad j=\lceil Kpc_P\rceil,
\]
where a fixed rational $K>C_1$ is chosen once.
Generic base change and shrinking make these sheaves vector bundles
of respective ranks
\begin{equation}\label{prod:image-ranks}
 \begin{split}
 &h^0(F,pD_F),\\
 &h^0(F,pD_F)-h^0(F,pD_F-jP).
 \end{split}
\end{equation}
In the second line we work on the fixed extraction.  The difference
comes from the kernel of restriction of the original section space;
it does not require surjectivity onto all sections of the thickening.

The test space injects into the jets of order $\lceil Kp\rceil$ of
this vector bundle at the chosen closed point of the base.  Indeed,
jets mean reduction modulo the $\lceil Kp\rceil$-th power of the
maximal ideal $\mathfrak m$ of the base point.  Vanishing in them,
followed by evaluation upstairs, puts the local section in
\[
 \mathfrak m^{\lceil Kp\rceil}\cO_{\mathcal X}(pD_W)
 +\mathcal I_{\mathcal P}^{\,j}\cO_{\mathcal X}(pD_W)
\]
in the divisor test, and in just the first summand in the component
test.  With the chosen weights both summands have order at least $Kp$, since
$j/c_P\geq Kp$.  This contradicts the bound $C_1p$.  Jets in $m$ smooth
base parameters have dimension at most $C_2p^m$ times the bundle rank,
with $C_2$ depending only on $m,K$.

Now $\dim R_{\leq ap}\sim hp^n/n!$.  Passing to volumes in the first
test gives the missing inequality in \eqref{prod:initial-volume};
when $q=0$, its rank is one and the same calculation applies.  In the
second test, for $q>0$, it gives
\begin{equation}\label{prod:volume-deficit}
 h\leq C\left[D_F^q-\vol(D_F-Kc_PP)\right].
\end{equation}
Here $P$ and its extraction are fixed before taking the large divisible
$p$ limit.  Rounding $Kpc_P$ does not change the leading term: one can
squeeze it between the corresponding rational coefficients differing
by an arbitrarily small positive amount and use continuity of volume.
The constants are uniform because they came from the jet estimate;
no uniform convergence in $P$ is required.

Combining \eqref{prod:volume-deficit} with the upper bound in
\eqref{prod:initial-volume} proves \eqref{prod:clipping}.  In detail,
if $t=P(J)>Kc_P$, set $\lambda=Kc_P/t$.  On the extraction,
\[
 D_F-Kc_PP\sim_{\Q}(1-\lambda)D_F+
                         \lambda J-Kc_PP,
\]
and the last two terms together form an effective divisor.
Monotonicity and homogeneity of volume imply
\[
 \vol(D_F-Kc_PP)\geq(1-\lambda)^qD_F^q.
\]
The two preceding bounds force
$1\leq C[1-(1-\lambda)^q]$, and hence force a uniform positive
lower bound on $\lambda$.  This is the claimed bound on $t/c_P$.
If $t\leq Kc_P$ it holds immediately.  There is no divisorial clipping
claim to check on a zero-dimensional component.

\subsection{Transfer across every later tier}

Suppose clipping has been proved on a prefix component $F$ of positive
dimension, and the next scale is $T$.  Set $L=TD_F$.  The bound
\eqref{prod:clipping} scales to hypothesis (iv) of
Proposition~\ref{prop:integral-jump}.  Its other hypotheses follow from
\eqref{prod:lc-poles}, Lemma~\ref{prod:one-section}, and
Lemma~\ref{prod:graph-bound}.  More explicitly, at the primes of the
normal graph we can take
\[
 \epsilon=\frac{\sum_{j\leq m}e_j-1}{T}\longrightarrow0,
\]
because every preceding degree is negligible compared with $T$.
All divisors and form supports were made SNC.  Therefore
\eqref{prod:integral-jump-bound} holds for $(F,\eta_F,L)$, after a
subsequence.

Cut by the $k_1$ pencils of the next tier, all with degrees at most
$CT$, and choose a geometric generic component $F^+$ of dimension
$q-k_1$.  Components and resolutions can be chosen compatibly by
taking successive algebraic closures of the rational base fields.
The next relative form is obtained by dividing by the new coordinate
differentials.  We claim
\begin{equation}\label{prod:tier-volume}
 L^q\leq C(L|_{F^+})^{q-k_1},\qquad
 (L|_{F^+})^{q-k_1}\leq CL^q,
\end{equation}
and, if $\dim F^+>0$, clipping for $L|_{F^+}$.

Apply the preceding jet argument on $F$ itself, testing all sections
of $pL$ and using only the $k_1$ new base parameters.  No torus
parameter is needed.  A nonzero section has effective divisor $pJ$
with $J\sim_{\Q}L$, so its local section order is bounded by
\[
 pC'\bigl(A_{\eta_F}(w)+w(L)\bigr)
\]
by \eqref{prod:integral-jump-bound}.  Spread any chosen geometric
generic divisor after finite base change and take sufficiently general
closed points avoiding ramification.  The estimates therefore remain
valid on this spreading family.  For the whole-component test the
weights one on the new base parameters give discrepancy $k_1$ and
divisor value zero.  At an lc place $P^+$ over $F^+$, give its equation
weight $1/P^+(L)$.  The relative form has a simple pole, so discrepancy
is again $k_1$, and now $w(L)=1$.

Use the pushforward bundle for the first test, and in the second use
the image of restriction to the thickening with order
$\lceil KpP^+(L)\rceil$.  The same weighted-order contradiction injects
the test space into jets in $k_1$ parameters.  Passing to volume gives
the first inequality of \eqref{prod:tier-volume} and the volume-deficit
estimate with source volume $L^q$.  Nef intersection with the new
pencils, each bounded by $CL$, gives the second inequality.  Combining
them with the effective-representative argument just given proves
clipping on $F^+$ for $L|_{F^+}=TD_{F^+}$, equivalently clipping for
$D_{F^+}$.  Higher resolutions do not alter these estimates: effective
representatives of a pulled-back $\Q$-linear system descend and pull
back exactly.  If $F^+$ is a point, the first test has rank and degree
one and gives the required volume comparison; the divisor test is
unnecessary.

This proves the transfer with constants uniform after subsequence.
There are only finitely many tiers, so finitely many subsequence choices
suffice.  To conclude, \eqref{prod:initial-volume} starts the induction
\[
 h\prod_{j\leq m}d_j\leq C D_F^{n-m},
\]
because the first-tier degrees are bounded.  Across a tier of length
$k_1$, the first inequality of \eqref{prod:tier-volume} says
\[
 T^qD_F^q\leq CT^{q-k_1}D_{F^+}^{q-k_1},
 \quad\text{hence}\quad
 T^{k_1}D_F^q\leq CD_{F^+}^{q-k_1}.
\]
Every new $d_j$ is comparable with $T$, giving exactly the new factors
in the product.  The terminal component is a point of degree one.
This proves Proposition~\ref{prop:product}.
